\begin{table}[htbp]
\centering
\begin{threeparttable}
\caption{Mean loss differences in the 2020 realised-information comparison.}
\label{tab:subsample-scores}
\small
\setlength{\tabcolsep}{3pt}
\begin{tabular*}{\textwidth}{@{\extracolsep{\fill}}llrrr@{}}
\toprule
Loss & Horizon & S\&P~500 & FTSE~100 & DAX \\
\midrule
@@ROW_0@@
@@ROW_1@@
@@ROW_2@@
@@ROW_3@@
@@ROW_4@@
@@ROW_5@@
@@ROW_6@@
\bottomrule
\end{tabular*}
\begin{tablenotes}[flushleft]
\footnotesize
\item Notes: Entries are mean losses for RV-NN-SV minus NN-SV, evaluated against the original Oxford--Man outcomes. Negative differences favour realised information. Brackets contain two-sided descriptive DM $p$-values, unadjusted for multiplicity. CRPS differences on the decimal-return scale are multiplied by $10^4$. Forecasts are assigned to 2020 by their origin dates. At horizons $(1,5,10)$, paired sample sizes are $(243,224,199)$ for the S\&P~500 and $(250,242,232)$ for both the FTSE~100 and DAX. All scores at a given market and horizon use the same paired sample.
\end{tablenotes}
\end{threeparttable}
\end{table}
